\bookchapter{Chapter 16: Account Systems and Cloud Services}

\chapterintro{Photos and videos}

\booksection{: From your child's birth to graduation ceremony, from first date to wedding trip.}

Notes

\begin{itemize}
  \item \textbf{Contacts}App purchase history
  \item \textbf{Photos and videos}Find My Phone
  \item \textbf{: After your phone is lost, this is the only tool that can help you locate, lock, and remotely wipe data.}If you use an Apple phone, all of the above is tied to your Apple ID. If you use Android, it's tied to your Huawei/Xiaomi/OPPO/vivo account, plus your Google account. These accounts are the "household register" of your digital life—lose the register, and many things become unprovable.
  \item \textbf{2. Three Must-Do Things, Right Now}Enable two-factor authentication
  \item \textbf{: Whether you use Apple or Android, the first thing to do is go to account settings and turn on "Two-Factor Authentication" or "Two-Step Verification". This is like adding a deadbolt to your door—password alone isn't enough, your phone also needs to receive a verification code. Even if someone guesses your password online, without your phone, they can't get in.}Bind a reliable recovery method
\end{itemize}

: Two-factor authentication is secure, but there's a risk—what if you lose your own phone? So the system will prompt you to set up a “recovery key” or “backup email/phone number”. This is your only way to get your account back in an emergency. Best to write it on paper, don't store it on your phone.

\booksection{Confirm “Find My Phone” is enabled}

\begin{enumerate}
  \item \textbf{: This feature can help you locate, lock, play sound, and even remotely wipe all data after your phone is lost. It's the last line of defense. Many people skip this step when buying a phone, and when they want to use it later, they find it wasn't enabled—and it's too late.}3. iCloud: The "Air" for Apple Users
  \item \textbf{If you use an iPhone, sooner or later you'll face a soul-searching question: what to do when iCloud storage is full?}Apple gives you 5GB for free, and honestly, this 5GB isn't even enough for one backup. So most people eventually receive that annoying popup: "iCloud Storage Full". At this point, you have three options:
  \item \textbf{Clean up}: Go to Settings Your Name iCloud Manage Storage to see who's taking up your space. Generally, photos and backups are the biggest. You can delete unnecessary old backups first.
\end{enumerate}

\booksection{Upgrade}

: 50GB is 6 yuan per month, 200GB is 21 yuan per month. Personally, I think 50GB is enough for most people, mainly for backups and photo thumbnails.

Workaround

\begin{itemize}
  \item \textbf{: If you don't want to pay, you can periodically export photos to your computer, then turn off iCloud photo sync. But then you lose the convenience of automatic photo backup.}An important reminder: iCloud Photo Library is not a "cloud drive"—it's "sync". If you delete a photo on your phone, it will also be deleted from iCloud. If what you want is "backup to cloud then clear phone", then iCloud Photos is not what you need—you need iCloud Backup, or a third-party cloud drive.
  \item \textbf{4. Don't Confuse Sync with Backup}These two words sound similar, but the consequences are completely different:
  \item \textbf{Sync}is bidirectional: Delete on phone, delete on cloud; delete on cloud, delete on phone. iCloud Photos, Contacts, Notes are all sync. The benefit is consistency across multiple devices; the downside is "once deleted, it's really gone".
\end{itemize}

Backup

\begin{tcbinfo}{iCloud+: The Privacy Trio for Paid Users}
As long as you upgrade to any paid iCloud tier (50GB/200GB/2TB, etc.), you automatically become an\textbf{iCloud+}user, unlocking three privacy protection features:

\begin{itemize}
  \item \textbf{Private Relay}: Apple's version of a "VPN." Safari browsing traffic is relayed through Apple servers before reaching the website, so no one (including Apple) can simultaneously see who you are and what you're accessing. Ideal for protecting privacy on public Wi-Fi. Settings path:\texttt{Settings > [Your Name] > iCloud > Private Relay}.
  \item \textbf{Hide My Email}: Generates one-time random email addresses when registering on websites/apps, with emails automatically forwarded to your real email. If you no longer want emails from a site, simply turn off that random address — your real email is never exposed. Settings path:\texttt{Settings > [Your Name] > iCloud > Hide My Email}.
  \item \textbf{Custom Email Domain}: You can use your own domain (e.g., @yourname.com) to send and receive iCloud email, suitable for personal branding or family email management.
\end{itemize}

Note: Private Relay is not available in some countries/regions, and domestic users may need alternative solutions.

\end{tcbinfo}

\begin{tcbwarn}{iCloud Advanced Data Protection (ADP) — True End-to-End Encryption}
In late 2022, Apple launched\textbf{iCloud Advanced Data Protection (ADP)}, extending end-to-end encryption (E2EE) from the original limited data like Health and Passwords to\textbf{the vast majority of iCloud data}— including Photos, Notes, iCloud Backup, device information, and 23 categories in total.

\begin{itemize}
  \item \textbf{What Is E2EE}: Encryption keys exist only on your devices — even Apple's servers don't have them. Even if Apple is hacked or law enforcement requests data, Apple can't produce decryptable plaintext data.
  \item \textbf{:}:\texttt{Settings > [Your Name] > iCloud > Advanced Data Protection}, and follow the prompts to set up Recovery Contacts/Recovery Key (\textbf{Must Do}— if you lose the key, your data can't be recovered either).
  \item \textbf{Notes}: (1) After enabling, you must properly safeguard the Recovery Key, otherwise data cannot be recovered if you lose account access; (2) Shared Albums, iCloud Shared Photo Library, and a few other features don't participate in E2EE; (3) Apple IDs in mainland China do not currently support ADP — you need to switch to an account in another region.
\end{itemize}

One-sentence recommendation: If you have important private data in your account and your account region supports it, strongly consider enabling ADP — this is the strongest cloud data protection available to ordinary users.

\end{tcbwarn}

\booksection{is one-way: The phone takes a “snapshot” of its current state and stores it in the cloud. After that, if you delete something on the phone, the copy in the cloud backup is still there. iCloud Backup, local computer backup—these are backups. The benefit is safer data; the downside is they take up space and can't update in real time.}

5. The Account Matrix for Android Users

\begin{itemize}
  \item \textbf{The situation on the Android side is more complicated than Apple, because your data may be scattered across several accounts:}Brand account
  \item \textbf{(Huawei/Xiaomi/OPPO/vivo): Manages your cloud space, Find My Phone, app store, wallet. Each brand has its own cloud service, usually a basic-free + paid-expansion model.}Google account
\end{itemize}

\booksection{: If you use Gmail, Google Photos, Google Contacts, these are all in your Google account. Google gives 15GB for free, more generous than Apple, but it's also shared with Gmail and Drive.}

WeChat/QQ

\begin{itemize}
  \item \textbf{: They are independent of the phone system. Your chat history and WeChat favorites are neither in the brand account nor in the Google account—they only exist on WeChat's own servers.}So when Android users do backups, they need to check separately: brand cloud backup, Google backup, WeChat chat history backup—all three are essential.
  \item \textbf{6. Family Sharing: Let Your Family Benefit Too}Apple's "Family Sharing" lets you share app purchases, music, movies, iCloud space with your family, and even share location. You can add up to 6 people. Setup method: Settings Your Name Family Sharing. Huawei and Xiaomi also have similar features, findable in account settings as "Family Sharing" or "Family Cloud".
  \item \textbf{A very practical scenario: Set up Family Sharing on your parents' phones, so you can always see their location (provided they agree). If an elderly person wanders off or can't be reached by phone, you have a location.}The Consequences of Account Theft Are More Serious Than You Think
\end{itemize}

If a hacker gets your account password, they can: remotely lock your phone and extort ransom, view your photos and contacts, make purchases with your payment methods. More troublingly, if your account is tied to too many things, recovering the account itself is a nightmare. So, two-factor authentication—go enable it now.

\booksection{Hands-on: Give Your Account a Comprehensive Health Check}

Open phone settings, go to the account page, confirm two-factor authentication is enabled.

Check whether “trusted phone numbers” and “recovery key” are set up; write the key on paper.

\begin{tcbwarn}{Go to iCloud or cloud space management, see what data is syncing, turn off what you don't need.}
Confirm “Find My Phone” is enabled.

\end{tcbwarn}

\begin{tcbpractice}{If you use Android, check whether brand cloud backup, Google backup, and WeChat backup are all covered.}
\begin{enumerate}
  \item Chapter 17: Linking Your Phone with Other Devices
  \item The phone in your hand is actually a super remote control and data center. It can not only make calls, but also connect your computer, TV, watch, speaker, and car together, so these devices are no longer islands. In this chapter, we'll talk about how to make them "team up".
  \item 1. Collaboration with Computer: Phone and Computer Are No Longer Two Worlds
  \item Honestly, what I used to hate most was taking a photo on my phone and wanting to edit it on the computer—I had to find a data cable and wait for transfer. Later I discovered that you don't actually need a cable at all; good integration is ten times faster than a cable.
  \item Apple ecosystem (AirDrop + Universal Clipboard)
\end{enumerate}

\end{tcbpractice}
